\hypertarget{cont11:paper-ix-analytic-regularity-of-the-local-parameter-solution}{%
\subsection{Chapter IX: analytic regularity of the local parameter
solution}\label{cont11:paper-ix-analytic-regularity-of-the-local-parameter-solution}}

The argument establishes both analyticity and differentiability of the
actual local polynomial solution map at a root with an invertible solved
Jacobian. The statements apply to every prime, including \(p=2\). They
also hold over any complete nontrivially normed field.

Let \(\mathbb K\) be such a field, let \(e,k\in\mathbb N\), and let \[
F\in\mathbb K[Y_1,\ldots,Y_e,T_1,\ldots,T_k]^e,
\qquad a=(y_0,t_0)\in\mathbb K^e\times\mathbb K^k.
\] Evaluation of \(F\) means coordinatewise evaluation of its actual
polynomials. Define the actual derivative blocks \[
J=D_YF(a)\in M_e(\mathbb K),
\qquad K=D_TF(a)\in M_{e\times k}(\mathbb K),
\] where each matrix entry is the corresponding formal partial
derivative evaluated at \(a\). The hypotheses are exactly \[
F(a)=0,
\qquad J\in\operatorname{GL}_e(\mathbb K).
\] Neither analyticity of a proposed solution nor a derivative formula
is assumed.

\textbf{Local analytic graph theorem.} There exist a map \[
\psi:\mathbb K^k\longrightarrow\mathbb K^e
\] and open sets \(U\subseteq\mathbb K^k\), \(V\subseteq\mathbb K^e\)
such that \[
t_0\in U,\qquad y_0\in V,\qquad \psi(t_0)=y_0,
\qquad \psi(U)\subseteq V.
\] The map is analytic at every point of \(U\), and its graph describes
all roots in the product neighborhood: \[
\forall t\in U\;\forall y\in V,
\qquad F(y,t)=0\quad\Longleftrightarrow\quad y=\psi(t).
\] At the base parameter it has the strict Fréchet derivative \[
\boxed{D\psi(t_0)=-J^{-1}K.}
\] Consequently it is differentiable on \(U\). The map is defined on the
ambient parameter space, but its root, analyticity and uniqueness
conclusions are restricted to the specified neighborhood. Uniqueness is
among solved coordinates in \(V\).

The main Lean theorem is
\mcode{Komorebi.PolynomialAnalyticImplicit.exists_analytic_implicit}.
The matrix in its \mcode{HasStrictFDerivAt} conclusion is the
continuous linear map represented by the displayed matrix.
\mcode{differentiableOn_of_analytic_graph} records the
differentiability consequence.

\textbf{Proof chain.} The scalar polynomial derivative bridge is proved
by induction on actual multivariate polynomials: constants, sums, and
multiplication by a variable. For a scalar polynomial \(P\), its strict
derivative is \[
v\longmapsto\sum_j \frac{\partial P}{\partial x_j}(a)v_j.
\] Assembling the coordinates gives the strict derivative of the
evaluated system as its actual Jacobian matrix. Polynomial analyticity
is supplied by the pinned Mathlib theorem for multivariate polynomial
evaluation, followed by finite coordinate assembly.

Consider the actual augmented polynomial map \[
\Phi(y,t)=(F(y,t),t).
\] Its derivative and inverse derivative at \(a\) are \[
D\Phi(a)=
\begin{pmatrix}J&K\\0&I_k\end{pmatrix},
\qquad
D\Phi(a)^{-1}=
\begin{pmatrix}J^{-1}&-J^{-1}K\\0&I_k\end{pmatrix}.
\] The proof constructs the continuous linear equivalence from these
explicit matrices and verifies both inverse identities. The strict
inverse theorem gives a local inverse \(g\) of \(\Phi\), and the
analytic inverse result shows that this same inverse is analytic near
\(\Phi(a)=(0,t_0)\). Set \[
\psi(t)=\pi_Y g(0,t).
\] The two inverse identities and a sufficiently small product
neighborhood give the exact local graph statement. Composing the
derivative of \(g\) with parameter insertion and solved-coordinate
projection yields \(-J^{-1}K\). The chosen parameter neighborhood is
also restricted so that the inverse identities, analyticity and
inclusion \(\psi(U)\subseteq V\) hold there simultaneously.

\textbf{Transfer to a supplied continuous solution.} Suppose
\(g:\mathbb K^k\to\mathbb K^e\) is continuous at \(t_0\),
\(g(t_0)=y_0\), and \[
F(g(t),t)=0
\] for all \(t\) in some neighborhood of \(t_0\). Under the same
invertibility hypothesis at \(a\), continuity places \(g(t)\) in the
uniqueness neighborhood \(V\) for all sufficiently close parameters.
Hence \(g\) and \(\psi\) agree on a neighborhood of \(t_0\). The theorem
\mcode{continuous_selection_regularity} transfers both analyticity
and the strict derivative formula to \(g\).

There is also an explicit open-set version. Let \(U\) be open, let \(g\)
be continuous on \(U\), and suppose for every \(t\in U\) that \[
F(g(t),t)=0,
\qquad
J_t:=D_YF(g(t),t)\in\operatorname{GL}_e(\mathbb K).
\] Writing \(K_t=D_TF(g(t),t)\), the theorem
\mcode{continuous_selection_analyticOnNhd} proves \[
g\text{ is analytic at every }t\in U,
\qquad Dg(t)=-J_t^{-1}K_t\quad(t\in U),
\] with a strict Fréchet derivative at each point. The invertibility hypothesis at each point is part of this theorem. This is a transfer
statement for continuous local root selections satisfying its
hypotheses. It does not assert equality of arbitrary global branch
choices.

\textbf{Integral p-adic specialization.} For any prime \(p\), let \[
F\in\mathbb Z_p[Y,T]^e,
\qquad a\in\mathbb Z_p^{e+k},
\qquad F(a)=0.
\] Let \(\iota:\mathbb Z_p\hookrightarrow\mathbb Q_p\) be the
coefficient inclusion, applied entrywise to points and matrices. Assume
only that \[
\iota(D_YF(a))\in\operatorname{GL}_e(\mathbb Q_p).
\] The theorem
\mcode{PolynomialRegularityIntegral.exists_analytic_integral_implicit}
constructs the preceding analytic graph in the ambient spaces
\(\mathbb Q_p^k\) and \(\mathbb Q_p^e\). Its root equation evaluates the
original integral polynomial coefficients using \(\iota\), and its base
derivative is \[
-\iota(D_YF(a))^{-1}\,\iota(D_TF(a)).
\] The source proves that polynomial evaluation, formal partial
differentiation and both Jacobian blocks commute with this inclusion. No
invertibility modulo \(p\) or integral inverse is required. This permits
a solved Jacobian that is singular modulo \(p\) while remaining
invertible over \(\mathbb Q_p\). The new specialization is an
ambient-field theorem. It does not treat \(\mathbb Z_p\) as a vector
space over \(\mathbb Q_p\) or assert a new explicit integral radius.

\begin{longtable}[]{@{}
  >{\raggedright\arraybackslash}p{(\columnwidth - 2\tabcolsep) * \real{0.5000}}
  >{\raggedright\arraybackslash}p{(\columnwidth - 2\tabcolsep) * \real{0.5000}}@{}}
\toprule\noalign{}
\begin{minipage}[b]{\linewidth}\raggedright
Source module
\end{minipage} & \begin{minipage}[b]{\linewidth}\raggedright
Mathematical role
\end{minipage} \\
\midrule\noalign{}
\endhead
\bottomrule\noalign{}
\endlastfoot
\mcode{PolynomialAnalyticEvaluation.lean} & Actual polynomial
analyticity and strict derivative equals the evaluated formal
Jacobian. \\
\mcode{PolynomialAnalyticInverse.lean} & Analyticity of the same local
inverse supplied by an invertible strict derivative. \\
\mcode{PolynomialAnalyticBlocks.lean} & Explicit augmented derivative
equivalence and its solved-coordinate restriction. \\
\mcode{PolynomialAnalyticImplicit.lean} & Open local root graph,
analyticity, base strict derivative and differentiability. \\
\mcode{PolynomialAnalyticSelection.lean} & Transfer to continuous root
selections, including the per-point open-set theorem. \\
\mcode{PolynomialRegularityIntegral.lean} & Integral coefficients and
centers included in the ambient p-adic field, with actual Jacobian
compatibility. \\
\end{longtable}

This block closes the local analytic regularity obligation stated in
\mcode{TARGET_CONTRACT.json}. It supplies no new uniform numerical
radius. The previous weighted lifting and joint deformation results
retain their own explicit radii and error estimates. It does not
establish a global analytic branch or a result at a point whose solved
Jacobian is singular over the fraction field. The proof uses established
analytic inverse machinery and makes no novelty claim for the general
implicit-function theorem.

The remaining Chapter IX work includes general singular persistence with
higher obstructions, native presentation invariance (13.2), global or
completed compatibility and descent (13.3), the outstanding sharper
arithmetic statements, and manuscript consolidation. None of those is
closed by this local regularity theorem. Release-level checks and audit
counts are recorded in the separate release verification and audit
files.
